\subsection{Game Object và cấu trúc dữ liệu bàn chơi}
\label{sec:game_objects}
\begin{sloppypar}

Module \texttt{GameObject/objects.py} khai báo năm lớp mô tả một đề bài. Tất cả
đều là \texttt{dataclass(frozen=True)}: không đối tượng nào sửa được sau khi
tạo, và các tập hợp bên trong được lưu bằng \texttt{tuple} thay vì
\texttt{list}. Nhờ đó một \texttt{PuzzleDefinition} được dùng chung an toàn
giữa nhiều ván, giữa engine và bộ kiểm tra, mà không cần sao chép. Toàn bộ
trạng thái thay đổi được của ván tách riêng sang \texttt{GameState}
(Mục~\ref{sec:game_state}).

\begin{table}[htbp]
  \centering\small
  \caption{Các lớp dữ liệu trong \texttt{objects.py}}
  \label{tab:game_objects}
  \begin{tabularx}{\textwidth}{l>{\raggedright\arraybackslash}X}
    \toprule
    \textbf{Lớp} & \textbf{Trường và ý nghĩa} \\
    \midrule
    \texttt{Cell} & \texttt{id: int}; \texttt{coordinates: (int, int)} là
      (hàng, cột) đánh số từ $0$; \texttt{character\_name}; \texttt{profession} \\
    \texttt{Board} & \texttt{cells}: bộ hai chiều các \texttt{Cell}, theo hàng.
      Thuộc tính tính toán \texttt{height}, \texttt{width}, \texttt{cell\_count} \\
    \texttt{Verdict} & \texttt{cell\_id}; \texttt{is\_criminal: bool}
      (\texttt{True} là tội phạm) \\
    \texttt{Clue} & \texttt{id}; \texttt{source\_cell\_id: int | None}
      (\texttt{None} là lời khai công khai từ đầu); \texttt{constraint};
      \texttt{description\_nl} là câu tiếng Việt hiển thị cho người chơi \\
    \texttt{PuzzleDefinition} & \texttt{id} (ví dụ \texttt{medium-4x4-harbor});
      \texttt{board}; \texttt{clues}; \texttt{solution} là đáp án đầy đủ;
      \texttt{initial\_verdicts} là các phán quyết có sẵn khi bắt đầu \\
    \bottomrule
  \end{tabularx}
\end{table}

\paragraph{Định danh và tọa độ.}
Mỗi ô có hai cách định vị: định danh phẳng \texttt{id} (khóa dùng chung giữa
engine, bộ kiểm tra và API) và \texttt{coordinates} (dùng khi cần quan hệ hình
học). \texttt{Board} chỉ lưu lưới ô và suy ra kích thước từ đó:
\texttt{height} là số hàng, \texttt{width} là độ dài hàng đầu, và
\texttt{cell\_count} $= \texttt{width}\times\texttt{height}$. Lớp này không tự
kiểm tra lưới có hình chữ nhật hay tọa độ khai báo có khớp vị trí. Hai điều kiện
đó được \texttt{validate\_puzzle} kiểm tra trong hàm \texttt{\_check\_grid}. Kiểm
tra này cần thiết vì \texttt{BruteForceChecker} và \texttt{CNFEncoder} tìm láng
giềng cạnh của \texttt{ConnectedConstraint} dựa trên \texttt{coordinates}, chứ
không dựa trên vị trí trong mảng.

\paragraph{Dựng bàn chơi cho các bản đồ.}
Sáu bản đồ dùng chung hàm \texttt{make\_board(size)} trong
\texttt{Maps/common.py}. Hàm lấy $\texttt{size}^2$ nhân vật đầu tiên từ danh
sách cố định $25$ cặp (tên, nghề), gán \texttt{id} tăng dần từ $0$ và tọa độ
\texttt{divmod(id, size)}, rồi cắt thành từng hàng. Vì vậy định danh luôn theo
thứ tự hàng và $\texttt{id} = r \cdot W + c$, đúng với công thức
\eqref{eq:cell_id}. Mỗi module bản đồ khai báo thêm tập \texttt{CRIMINALS},
các \texttt{Clue}, và dựng \texttt{solution} bằng
\texttt{Verdict(i, i in CRIMINALS)} cho mọi ô.

\paragraph{Truy vấn trên bàn chơi.}
Module \texttt{board\_utils.py} cung cấp bốn hàm trả về bộ định danh ô:
\texttt{row\_ids(board, r)}, \texttt{column\_ids(board, c)},
\texttt{profession\_ids(board, p)}, và \texttt{neighbor\_ids(board, i)} trả về
tối đa tám ô lân cận kể cả đường chéo, không gồm chính ô $i$. Hàm cuối duyệt
lưới tìm ô $i$, rồi cắt lát các hàng và cột trong khoảng
$[\max(0, r-1),\, r+2)$ và $[\max(0, c-1),\, c+2)$, nên tự xử lý ô ở biên và
góc. Ô không tồn tại gây \texttt{ValueError}. Các hàm này mô tả đúng những tập
đích mà lời khai trong bản đồ dùng tới (hàng, cột, lân cận). Tuy nhiên, ở phiên
bản hiện tại các module bản đồ ghi tường minh tập định danh trong từng ràng buộc,
ví dụ \texttt{Count((3, 4, 5))} cho hàng 2 của lưới $3\times3$, và không module
nào gọi tới \texttt{board\_utils}.
\end{sloppypar}
